\documentclass{situationist-preprint}

\title{Why the Traveller Fails}
\msid{TSM-04-b}
\msfamily{TSM-04}
\msdiscipline{Time, Space and Modality}
\mstype{Companion manuscript}
\mskeywords{time travel; grandfather paradox; coincidence; Horwich; explanation; selection}

\begin{document}
\maketitle

\begin{abstract}
TSM-04-a proves that in every consistent history in which a time traveller arrives, her grandfather survives, so that if she tries, circumstances defeat her. Critics hold that this makes time travel depend on inexplicable coincidences: banana peels, jammed guns and sudden changes of heart. We argue that each failure has an ordinary local cause, that the pattern of failures is explained by the traveller's arrival, and that what remains improbable is the occurrence of certain loops, not the failures within them. The improbability bears on how often time travel happens, not on whether it is possible.
\end{abstract}

\section{The coincidence worry}

In the model of TSM-04-a, the setting in which the traveller arrives, tries, finds favorable circumstances and later departs has no consistent history (Proposition~3.4 there), and in every consistent history with her arrival, the circumstances are unfavorable if she tries (Theorem~3.3). \citet{Horwich1987} argues that this makes time travel improbable: if a traveller tried many times, she would need to fail each time, by a sequence of independent mishaps whose conjunction has negligible probability. \citet{Smith1997} replies that the improbability is an artefact of the choice of reference class.

\section{The failures are explained}

\begin{proposition}\label{p:local}
Each failure has a cause that is local and of an ordinary kind.
\end{proposition}

\begin{proof}
In a consistent history, the traveller's attempt is defeated by some event in its vicinity: the gun misfires, she is distracted, she mistakes another man for her grandfather. Each such event has its ordinary antecedents, and the consistency condition of TSM-04-a (Principle~P1 there) imposes no event that lacks them; it selects which histories there are, and each of them is lawful throughout.
\end{proof}

\begin{proposition}\label{p:pattern}
The pattern of failures is explained by her arrival.
\end{proposition}

\begin{proof}
Her arrival entails her later birth and departure, hence her grandfather's survival (Theorem~3.3 of TSM-04-a). Conditional on her arrival, the probability that every attempt fails is one. The correlation among the failures is the correlation that any set of events bears to a condition they are jointly required to satisfy, and the condition is part of the history.
\end{proof}

Propositions~\ref{p:local} and~\ref{p:pattern} together answer the question of the title. She fails because something ordinary stops her on each occasion, and she fails on every occasion because she arrived. The second explanation is not a further mechanism acting on the first. It is the same kind of explanation as that of why a randomly chosen survivor of a lottery-like process has the traits survival requires: by selection, not by an extra force.

\section{Objection and reply}

\begin{objection}
Selection explains why, given the loop, she fails; it does not remove the improbability of a loop containing many attempts, which on the dynamics alone would rarely end in failure every time.
\end{objection}

\begin{reply}
Granted. The unconditional probability of a loop with many attempts may be very small, and this is Horwich's point in its correct form. It bears on how often such loops occur, and perhaps on the evidence that any traveller would need before believing she had arrived. It does not bear on their possibility, and TSM-04 claims only that backward time travel is consistent. Nor does it show that the traveller cannot try: as \citet{Lewis1976} argues, relative to her skill and equipment she can kill him, and relative to the fact of her arrival she cannot \citep[cf.][]{Vihvelin1996}.
\end{reply}

\section{Conclusion}

The traveller's failures are lawful and locally caused, and their pattern is explained by her arrival. What is improbable is the loop as a whole, which limits how often time travel occurs, not whether it can.

\begin{thebibliography}{9}
\bibitem[Horwich(1987)]{Horwich1987} Horwich, P. (1987). \emph{Asymmetries in Time}. MIT Press.
\bibitem[Lewis(1976)]{Lewis1976} Lewis, D. (1976). The paradoxes of time travel. \emph{American Philosophical Quarterly}, 13(2), 145--152.
\bibitem[Smith(1997)]{Smith1997} Smith, N.~J.~J. (1997). Bananas enough for time travel? \emph{The British Journal for the Philosophy of Science}, 48(3), 363--389.
\bibitem[Vihvelin(1996)]{Vihvelin1996} Vihvelin, K. (1996). What time travelers cannot do. \emph{Philosophical Studies}, 81(2--3), 315--330.
\end{thebibliography}

\end{document}
